\documentclass[11pt,a4paper]{article}
\usepackage[utf8]{inputenc}
\usepackage[T1]{fontenc}
\usepackage{amsfonts}
\usepackage[french]{babel}
\frenchbsetup{StandardLists=true} % à inclure si on utilise \usepackage[francais]{babel}
\usepackage{enumitem}
\usepackage{amssymb}
\usepackage{amsmath}
\usepackage[version=4]{mhchem}
\usepackage[left=2cm,right=2cm,top=2cm,bottom=2cm]{geometry}
\usepackage{multirow}
\usepackage[dvipsnames]{xcolor}

\usepackage[T1]{fontenc}
\usepackage{fouriernc}
\usepackage{graphicx}

\usepackage{setspace}
\singlespacing

\usepackage{pstricks}
\usepackage{xkeyval}
\usepackage{pst-labo}


\usepackage{multicol}
\usepackage{caption}
\usepackage{fancyhdr}
\pagestyle{fancy}
\renewcommand\headrulewidth{1pt}
\fancyhead[L]{Lycée Denis Diderot }
\fancyhead[R]{Classe de Terminale}
\setcounter{page}{3}\fancyfoot[C]{page \thepage}
\fancyfoot[R]{Année 2025 2026}
\fancyfoot[L]{Alexis RUIZ}
\usepackage{multirow,tabularx,booktabs}
\newcommand{\cadreblanc}[1]{%
\medskip\framebox[\linewidth][c]{%
\rule{0mm}{#1mm}}\par
\medskip}

\usepackage{awesomebox}
\setlength{\aweboxvskip}{0mm}
\usepackage{pifont}
\usepackage{chemmacros}
\chemsetup{modules=all}

\renewcommand{\arraystretch}{2}
\newcommand*\acocher{\quitvmode{\fboxsep0pt \fboxrule0.8pt \fbox{\vrule height1.8ex depth.1ex width0pt \vrule height0pt depth0pt width1.9ex }}}

\newcommand{\defproblem}{\awesomebox[black]{2pt}{\rotatebox{90}{\large{\textbf{Problématique}}}}{black}}
\newcommand{\defconclusion}{\awesomebox[black]{2pt}{\rotatebox{90}{\Large{\textbf{Conclusion}}}}{black}}

\usepackage{tcolorbox}

\newcommand{\cor}[1]{%
  \begin{tcolorbox}[
    colback=red!10,
    colframe=red,
    arc=3mm,
    boxrule=1pt,
    left=5pt,
    right=5pt,
    top=15pt,
    bottom=5pt
  ]
  #1
  \end{tcolorbox}
}

\frenchbsetup{StandardLists=true}
\usepackage{titlesec}
\renewcommand{\thesection}{\ding{\numexpr201+\value{section}\relax}}

\titleformat{\section}
  {\normalfont\large\bfseries}
  {\thesection}
  {0.2em}
  {}
  
\begin{document}


\center \LARGE{COURS - Prévoir une réaction d'oxydo-réduction}
\\[0,3cm]

\normalsize

\hrule\
\\[1cm]

\flushleft
\section{Corrosion des métaux : une réaction d'oxydoréduction}

Lors du TP nous avons vu : \\[0.8cm]

\textcolor{red}{Il y a transfert d'électrons entre le métal fer et les ions \ce{Cu^{2+}}, selon les demi-équations électroniques :}  \\[0.8cm]

\cor{
\vspace{-8mm}
\Large
\begin{align*}
\ce{Fe &-> Fe^{2+} + 2 e^-} \quad &\text{(oxydation)} \\[3mm]
\ce{Cu^{2+} + 2 e^- &-> Cu} \quad &\text{(réduction)}
\end{align*}
}

\vspace{5mm}

\normalsize
\begin{itemize}
\item Ces 2 phénomènes sont simultanés et spontanés. Ils se traduisent par l'équation-bilan : \\[0.5cm]

\Large{\textcolor{red}{
\begin{equation*}
\ce{Fe + Cu^{2+} -> Fe^{2+} + Cu} \quad \text{(oxydo-réduction)}
\end{equation*}
}
}\end{itemize}


\section{Définitions}

\cor{

\begin{itemize}[leftmargin=1cm]
\item L'oxydation est une perte d'électrons

\vspace{8mm}

\item La réduction est un gain d'électrons

\vspace{8mm}

\item L'oxydant est l'espèce chimique (atome, ion ou molécule) qui gagne des électrons

\vspace{8mm}

\item Le réducteur est l'espèce chimique qui donne des électrons

\vspace{8mm}

\item L'ensemble des formes oxydées et réduites constitue un couple oxydo-réducteur \ce{M^{n+}/M} appelé couple rédox.

\vspace{8mm}

\item Par convention, l'oxydant figure en premier dans l'écriture du couple rédox 

\vspace{8mm}

(exemple : \ce{Cu^{2+}/Cu} ; \ce{Ag^+/Ag} ; \ce{Fe^{2+}/Fe})
\end{itemize}



}

\end{document}